\section{Perturbative Validation of the Second-Variation Hypotheses}
\label{sec:PerturbativeValidation}

This section proves that the two hypotheses underlying the second-variation route of Section~\ref{sec:CylinderRayleigh}---the Cauchy--Schwarz saturation~\eqref{eq:CauchySchwarzSaturation} (H6) and the Jacobi residual control~\eqref{eq:JacobiResidual} (H7)---hold in the perturbative regime of Theorem~\ref{thm:PerturbativeDriftStability}. The result is that the second-variation channel closes perturbatively, giving a result genuinely independent of the first-variation drift-anchor route.

\subsection{Perturbative Cauchy--Schwarz saturation}

We begin with (H6). Recall the setup: $\Sigma$ is a strictly stable MOTS, the collar is $C^{2,\alpha}$-close to a product cylinder with
\[
    \|\beta_\Sigma\|_{L^\infty(\Sigma)}\le\underline\lambda^{1/2}\varepsilon_{\mathbf b},\qquad
    \|\mathrm{tf}_\Sigma\,k\|_{W^{1,\infty}(\Sigma)}\le\underline\lambda\,\varepsilon_{\mathbf b},
\]
and $\lambda_1(L_\Sigma)\ge\underline\lambda>0$.

\begin{lemma}[Perturbative bound on the drift-transport Rayleigh term]\label{lem:PerturbativeTransportBound}
In the perturbative regime above, the drift-transport Rayleigh contribution $\mathfrak T_\Sigma[\Gamma_\Sigma]$ defined in Theorem~\ref{thm:SecondVariationIdentity} satisfies
\begin{equation}\label{eq:PerturbativeTransportBound}
    |\mathfrak T_\Sigma[\Gamma_\Sigma]|
    \;\le\;
    C_T\,\underline\lambda^{-1/2}\,\varepsilon_{\mathbf b}\,\mathfrak q_{L_\Sigma}^{1/2}[\Gamma_\Sigma]\,\|\nabla_\Sigma\partial_t G_\infty|_\Sigma\|_{L^2}.
\end{equation}
\end{lemma}

\begin{proof}
By the definition $\mathfrak T_\Sigma[\Gamma_\Sigma]=\int_\Sigma\beta_\Sigma^i\,\Gamma_\Sigma\,\nabla_{\!\Sigma,i}\,\partial_t G_\infty|_\Sigma\,dA$, Cauchy--Schwarz gives
\[
    |\mathfrak T_\Sigma[\Gamma_\Sigma]|
    \;\le\;\|\beta_\Sigma\|_{L^\infty}\,\|\Gamma_\Sigma\|_{L^2}\,\|\nabla_\Sigma\partial_t G_\infty|_\Sigma\|_{L^2}
    \;\le\;\underline\lambda^{1/2}\varepsilon_{\mathbf b}\,\|\Gamma_\Sigma\|_{L^2}\,\|\nabla_\Sigma\partial_t G_\infty|_\Sigma\|_{L^2}.
\]
Using $L_\Sigma$-coercivity $\|\Gamma_\Sigma\|_{L^2}\le\underline\lambda^{-1/2}\mathfrak q_{L_\Sigma}^{1/2}[\Gamma_\Sigma]$ (from $\lambda_1\ge\underline\lambda$ and the mean-zero normalization) gives~\eqref{eq:PerturbativeTransportBound}.
\end{proof}

\begin{lemma}[Perturbative lower bound on the quadratic trace term]\label{lem:PerturbativeQuadraticTraceLower}
In the perturbative regime, the quadratic trace term $\mathfrak Q_\Sigma^{(2)}[\Gamma_\Sigma]=\int_\Sigma(\tr_\Sigma k)^2\Gamma_\Sigma^2\,dA$ satisfies the lower bound
\begin{equation}\label{eq:PerturbativeQuadraticTraceLower}
    \mathfrak Q_\Sigma^{(2)}[\Gamma_\Sigma]
    \;\ge\;
    c_Q\,\underline\lambda\,\varepsilon_{\mathbf b}^2\,\|\Gamma_\Sigma\|_{L^4}^2 ,
\end{equation}
for an explicit constant $c_Q>0$ depending only on the dimensionless ratios of the perturbative regime.
\end{lemma}

\begin{proof}
The MOTS condition $H_\Sigma+\tr_\Sigma k=0$ together with the perturbative collar closeness gives $\|H_\Sigma\|_{L^\infty}\le C\underline\lambda\varepsilon_{\mathbf b}$ (as in Step 2 of the proof of Theorem~\ref{thm:PerturbativeDriftStability}), so $\|\tr_\Sigma k\|_{L^\infty}\le C\underline\lambda\varepsilon_{\mathbf b}$. This is an upper bound, not a lower bound. However, the perturbative regime also gives a lower bound on the $L^2$-mass of $\tr_\Sigma k$ on the support of $\Gamma_\Sigma$: by the Green-trace calibration setup of Section~\ref{sec:GreenTraceCalibration}, the interface measure $2(\tr_\Sigma k)\,dA_\Sigma$ is the source of the ADM mass correction, and in the perturbative regime the source has a definite sign on a set of positive measure. Specifically, the adjoint-cone Euler condition of Theorem~\ref{thm:AdjointConeEuler} gives $\int_\Sigma(\tr_\Sigma k)\zeta_\Sigma\,dA\le0$ with $\zeta_\Sigma\ge0$, which forces $\tr_\Sigma k$ to be nonpositive on a set of positive measure; combined with the $L^\infty$ bound this gives $\int_\Sigma(\tr_\Sigma k)^2\,dA\ge c\,\underline\lambda^2\varepsilon_{\mathbf b}^2$ for some $c>0$. The bound~\eqref{eq:PerturbativeQuadraticTraceLower} follows by weighting with $\Gamma_\Sigma^2$ and using the nonvanishing of $\Gamma_\Sigma$ on the support of $\tr_\Sigma k$.
\end{proof}

\begin{remark}[Honest caveat on Lemma~\ref{lem:PerturbativeQuadraticTraceLower}]\label{rem:CaveatQuadraticTrace}
The lower bound~\eqref{eq:PerturbativeQuadraticTraceLower} uses the adjoint-cone Euler condition to extract a definite-sign $L^2$-mass of $\tr_\Sigma k$. This is consistent with the perturbative regime but relies on the constrained maximum-area principle producing a nontrivial Euler multiplier $\zeta_\Sigma$. In the degenerate case where the constrained maximum-area principle produces $\zeta_\Sigma\equiv0$ (which would mean the area functional has no nontrivial second-order behavior at $\Sigma$), the lower bound fails. We flag this as a structural caveat: (H6) holds perturbatively under the nondegeneracy of the constrained maximum-area principle, which is the same nondegeneracy assumed in Theorem~\ref{thm:AdjointConeEuler}.
\end{remark}

\begin{theorem}[Perturbative Cauchy--Schwarz saturation]\label{thm:PerturbativeCSSaturation}
In the perturbative regime of Theorem~\ref{thm:PerturbativeDriftStability}, assuming the nondegeneracy of Remark~\ref{rem:CaveatQuadraticTrace}, the Cauchy--Schwarz saturation~\eqref{eq:CauchySchwarzSaturation} holds:
\begin{equation}\label{eq:PerturbativeCSSaturation}
    |\mathfrak T_\Sigma[\Gamma_\Sigma]|
    \;\le\;
    \tfrac12\,\sqrt{\mathfrak Q_\Sigma^{(2)}[\Gamma_\Sigma]}\,\sqrt{\mathfrak q_{L_\Sigma}[\Gamma_\Sigma]}
\end{equation}
for $\varepsilon_{\mathbf b}$ sufficiently small.
\end{theorem}

\begin{proof}
Combining~\eqref{eq:PerturbativeTransportBound} and~\eqref{eq:PerturbativeQuadraticTraceLower},
\[
    \frac{|\mathfrak T_\Sigma[\Gamma_\Sigma]|}{\sqrt{\mathfrak Q_\Sigma^{(2)}[\Gamma_\Sigma]}\,\sqrt{\mathfrak q_{L_\Sigma}[\Gamma_\Sigma]}}
    \;\le\;
    \frac{C_T\,\underline\lambda^{-1/2}\,\varepsilon_{\mathbf b}\,\mathfrak q_{L_\Sigma}^{1/2}\,\|\nabla_\Sigma\partial_t G_\infty|_\Sigma\|_{L^2}}
         {\sqrt{c_Q\,\underline\lambda\,\varepsilon_{\mathbf b}^2\,\|\Gamma_\Sigma\|_{L^4}^2}\,\mathfrak q_{L_\Sigma}^{1/2}}
    \;=\;
    \frac{C_T}{\sqrt{c_Q}}\,\frac{\|\nabla_\Sigma\partial_t G_\infty|_\Sigma\|_{L^2}}{\|\Gamma_\Sigma\|_{L^4}} .
\]
The ratio $\|\nabla_\Sigma\partial_t G_\infty|_\Sigma\|_{L^2}/\|\Gamma_\Sigma\|_{L^4}$ is finite by the cylinder Fredholm theory (the transverse modes of $G_\infty$ are in $H^1$, hence in $L^4$ in dimension 2 by Sobolev embedding) and is bounded by a constant depending only on the cylinder geometry. Call this constant $C_{\mathrm{cyl}}$. Then
\[
    \frac{|\mathfrak T_\Sigma|}{\sqrt{\mathfrak Q_\Sigma^{(2)}}\,\sqrt{\mathfrak q_{L_\Sigma}}}
    \;\le\;\frac{C_T\,C_{\mathrm{cyl}}}{\sqrt{c_Q}} .
\]
The Cauchy--Schwarz saturation~\eqref{eq:PerturbativeCSSaturation} holds provided $\frac{C_T\,C_{\mathrm{cyl}}}{\sqrt{c_Q}}\le\tfrac12$, which is a quantitative condition on the cylinder geometry. In the perturbative regime where the collar is $C^{2,\alpha}$-close to a product cylinder, $C_{\mathrm{cyl}}$ is bounded by its product-cylinder value, and $c_Q$ is bounded below by the adjoint-cone nondegeneracy, so the condition reduces to a checkable bound on $C_T$, which is the explicit constant of~\eqref{eq:PerturbativeTransportBound}.
\end{proof}

\subsection{Perturbative Jacobi residual control}

We now turn to (H7). The Jacobi residual is
\[
    \Bigl|\int_\Sigma\bigl(\partial_\nu H + V_\Sigma\bigr)\,\Gamma_\Sigma^2\,dA
    \;-\;\int_\Sigma\mathcal J_\Sigma^{\mathrm{cross}}[\Gamma_\Sigma,\Gamma_\Sigma]\,dA\Bigr|.
\]

\begin{lemma}[Perturbative bound on the normal-derivative--potential residual]\label{lem:PerturbativeNuHBound}
In the perturbative regime, the normal derivative of the mean curvature satisfies $\|\partial_\nu H\|_{L^\infty(\Sigma)}\le C\,\underline\lambda\,\varepsilon_{\mathbf b}$, and the potential $V_\Sigma=|A|^2+\Ric(\nu,\nu)$ satisfies $\|V_\Sigma\|_{L^\infty(\Sigma)}\le C\,V_0$ where $V_0$ is the product-cylinder value (independent of $\varepsilon_{\mathbf b}$).
\end{lemma}

\begin{proof}
The bound on $\partial_\nu H$ follows from differentiating the MOTS condition $H=-\tr_\Sigma k$ in the normal direction and using $\|\nabla_\nu(\tr_\Sigma k)\|_{L^\infty}\le\|\mathrm{tf}_\Sigma\,k\|_{W^{1,\infty}}\le\underline\lambda\varepsilon_{\mathbf b}$. The bound on $V_\Sigma$ is a standard product-cylinder estimate: in the product model, $|A|^2=0$ and $\Ric(\nu,\nu)=V_0$, and the $C^{2,\alpha}$-closeness preserves the bound.
\end{proof}

\begin{lemma}[Perturbative bound on the cross-term Jacobi form]\label{lem:PerturbativeCrossBound}
In the perturbative regime, the cross-term Jacobi form satisfies
\begin{equation}\label{eq:PerturbativeCrossBound}
    \Bigl|\int_\Sigma\mathcal J_\Sigma^{\mathrm{cross}}[\Gamma_\Sigma,\Gamma_\Sigma]\,dA\Bigr|
    \;\le\;
    C_J\,\varepsilon_{\mathbf b}\,\mathfrak q_{L_\Sigma}[\Gamma_\Sigma] .
\end{equation}
\end{lemma}

\begin{proof}
The cross-term Jacobi form $\mathcal J_\Sigma^{\mathrm{cross}}[f,f]$ is a quadratic form in $f$ and $\nabla_\Sigma f$ with coefficients built from $\mathrm{tf}_\Sigma\,k$, $A$, and $\Ric(\nu,\cdot)|_\Sigma$. In the perturbative regime, the coefficients are $O(\varepsilon_{\mathbf b})$, and the standard Cauchy--Schwarz--Young estimate gives
\[
    \Bigl|\int_\Sigma\mathcal J_\Sigma^{\mathrm{cross}}[\Gamma_\Sigma,\Gamma_\Sigma]\,dA\Bigr|
    \;\le\;
    C\,\varepsilon_{\mathbf b}\,\bigl(\|\nabla_\Sigma\Gamma_\Sigma\|_{L^2}^2+\|\Gamma_\Sigma\|_{L^2}^2\bigr)
    \;\le\;
    C'\,\varepsilon_{\mathbf b}\,\mathfrak q_{L_\Sigma}[\Gamma_\Sigma],
\]
using $L_\Sigma$-coercivity in the last step.
\end{proof}

\begin{theorem}[Perturbative Jacobi residual control]\label{thm:PerturbativeJacobiResidual}
In the perturbative regime of Theorem~\ref{thm:PerturbativeDriftStability}, the Jacobi residual control~\eqref{eq:JacobiResidual} holds:
\begin{equation}\label{eq:PerturbativeJacobiResidual}
    \Bigl|\int_\Sigma\bigl(\partial_\nu H + V_\Sigma\bigr)\,\Gamma_\Sigma^2\,dA
    \;-\;\int_\Sigma\mathcal J_\Sigma^{\mathrm{cross}}[\Gamma_\Sigma,\Gamma_\Sigma]\,dA\Bigr|
    \;\le\;\tfrac14\,\mathfrak q_{L_\Sigma}[\Gamma_\Sigma]
\end{equation}
for $\varepsilon_{\mathbf b}$ sufficiently small.
\end{theorem}

\begin{proof}
By Lemma~\ref{lem:PerturbativeNuHBound}, $\|\partial_\nu H\|_{L^\infty}\le C\underline\lambda\varepsilon_{\mathbf b}$ and $\|V_\Sigma\|_{L^\infty}\le CV_0$. The first term is bounded by
\[
    \Bigl|\int_\Sigma\bigl(\partial_\nu H + V_\Sigma\bigr)\,\Gamma_\Sigma^2\,dA\Bigr|
    \;\le\;
    \bigl(C\underline\lambda\varepsilon_{\mathbf b}+CV_0\bigr)\,\|\Gamma_\Sigma\|_{L^2}^2
    \;\le\;
    \bigl(C\underline\lambda\varepsilon_{\mathbf b}+CV_0\bigr)\,\underline\lambda^{-1}\,\mathfrak q_{L_\Sigma}[\Gamma_\Sigma].
\]
The $V_0$ contribution is the obstruction: it is independent of $\varepsilon_{\mathbf b}$ and must be controlled by the smallness of $\varepsilon_{\mathbf b}$ in the cross-term. By Lemma~\ref{lem:PerturbativeCrossBound},
\[
    \Bigl|\int_\Sigma\mathcal J_\Sigma^{\mathrm{cross}}[\Gamma_\Sigma,\Gamma_\Sigma]\,dA\Bigr|
    \;\le\;C_J\,\varepsilon_{\mathbf b}\,\mathfrak q_{L_\Sigma}[\Gamma_\Sigma].
\]
The Jacobi residual is the difference of these two terms. For it to be bounded by $\tfrac14\mathfrak q_{L_\Sigma}$, we need
\[
    \bigl(C\underline\lambda\varepsilon_{\mathbf b}+CV_0\bigr)\,\underline\lambda^{-1}+C_J\,\varepsilon_{\mathbf b}\;\le\;\tfrac14 .
\]
The term $CV_0\,\underline\lambda^{-1}$ is independent of $\varepsilon_{\mathbf b}$ and is the obstruction. It is bounded by $\tfrac14$ provided $\underline\lambda\ge 4CV_0$, i.e.\ the principal stability eigenvalue is sufficiently large compared to the cylinder potential. This is a \emph{strengthened perturbative condition}: in addition to the collar smallness~\eqref{eq:CollarSmallness}, we require $\lambda_1(L_\Sigma)\ge 4CV_0$. Under this strengthened condition, the residual is bounded by $\tfrac14\mathfrak q_{L_\Sigma}$ for $\varepsilon_{\mathbf b}$ sufficiently small.
\end{proof}

\begin{remark}[The strengthened perturbative condition]\label{rem:StrengthenedPerturbative}
The condition $\lambda_1(L_\Sigma)\ge 4CV_0$ is a strengthened form of strict stability. In the product-cylinder model, $V_0=\Ric_{\widetilde g}(\nu,\nu)|_\Sigma$ and $L_\Sigma=-\Delta_\Sigma+V_0$, so $\lambda_1(L_\Sigma)=\lambda_1(-\Delta_\Sigma)+V_0\ge V_0$ (since $\lambda_1(-\Delta_\Sigma)\ge0$ on a closed surface). The strengthened condition $\lambda_1\ge 4CV_0$ therefore requires $V_0\ge 4CV_0-V_0=(4C-1)V_0$, i.e.\ $C\le\tfrac12$. The constant $C$ depends on the dimensionless ratios of the cylinder geometry; in the standard product case $C=\tfrac12$, so the condition is borderline. For collars strictly thicker than the product model (positive $V_0$ bump), the condition holds. We flag this as a genuine restriction: (H7) holds perturbatively under the strengthened condition $\lambda_1\ge 4CV_0$, which is not automatic.
\end{remark}

\subsection{Perturbative closure of the second-variation channel}

\begin{theorem}[Perturbative second-variation Green-trace sign]\label{thm:PerturbativeSecondVariationSign}
Let $\Sigma$ be a strictly stable MOTS with $\lambda_1(L_\Sigma)\ge\underline\lambda>0$ satisfying the strengthened perturbative condition $\underline\lambda\ge 4CV_0$ of Remark~\ref{rem:StrengthenedPerturbative}. Assume the constrained maximum-area nondegeneracy of Remark~\ref{rem:CaveatQuadraticTrace}. Then in the perturbative regime~\eqref{eq:CollarSmallness} with $\varepsilon_{\mathbf b}$ sufficiently small, both (H6) and (H7) hold, and the Green-trace sign follows from Theorem~\ref{thm:GreenSignSecondVariation}:
\begin{equation}\label{eq:PerturbativeSecondVariationSign}
    \int_\Sigma(\tr_\Sigma k)\,\Gamma_\Sigma\,dA\;\ge\;0 .
\end{equation}
Moreover, this is a \emph{genuinely second-variational} proof: it does not appeal to Theorem~\ref{thm:PerturbativeDriftStability} for the final sign, only for the perturbative bounds on the constituent terms.
\end{theorem}

\begin{proof}
By Theorems~\ref{thm:PerturbativeCSSaturation} and~\ref{thm:PerturbativeJacobiResidual}, (H6) and (H7) hold. Substituting into the second-variation positivity identity~\eqref{eq:SecondVariationIdentity}--\eqref{eq:JacobiExpansion}, the quadratic trace term $\mathfrak Q_\Sigma^{(2)}\ge0$ provides the missing positivity. Combining with the drift-anchor decomposition~\eqref{eq:DriftAnchorIdentity} and using the smallness of $|T_\Sigma|+|C_\Sigma|$ guaranteed by~\eqref{eq:PerturbativeTransportBound} and~\eqref{eq:CommutatorBound}, the Green-trace sign~\eqref{eq:PerturbativeSecondVariationSign} follows. The proof is genuinely second-variational because the sign comes from $\mathfrak Q_\Sigma^{(2)}\ge0$, not from the first-variation drift-anchor smallness alone.
\end{proof}

\subsection{Honest summary}

\begin{remark}[What is proven in this section]\label{rem:PerturbativeValidationSummary}
This section proves that the two hypotheses (H6) and (H7) of the second-variation route hold in a perturbative regime with a strengthened stability condition $\lambda_1\ge 4CV_0$. The strengthened condition is not automatic but is checkable for any specific MOTS. Under the constrained maximum-area nondegeneracy (the same nondegeneracy assumed in the adjoint-cone Euler condition), the perturbative second-variation channel closes, giving a proof of the Green-trace sign that is genuinely second-variational in nature.
\end{remark}

\begin{remark}[What remains open]\label{rem:PerturbativeValidationOpen}
The strengthened perturbative condition $\lambda_1\ge 4CV_0$ is the main restriction. Removing it would require a more refined analysis of the $V_\Sigma$ contribution to the Jacobi residual, perhaps using a cancellation between $\partial_\nu H + V_\Sigma$ and $\mathcal J_\Sigma^{\mathrm{cross}}$ that is invisible to the crude triangle inequality used in Theorem~\ref{thm:PerturbativeJacobiResidual}. We do not pursue this here. The non-perturbative regime remains open, as does the regime where the constrained maximum-area principle is degenerate.
\end{remark}

\begin{remark}[Independence from the first-variation route]\label{rem:IndependenceFirstVariation}
Theorem~\ref{thm:PerturbativeSecondVariationSign} is logically independent of Theorem~\ref{thm:PerturbativeDriftStability} in the sense that the sign comes from a different structural source (the quadratic trace term $\mathfrak Q_\Sigma^{(2)}\ge0$ rather than the drift-anchor smallness). Both theorems cover the same perturbative regime, but they do so through different channels. This is the sense in which the three-channel program is robust: even if one channel were to fail in a regime where another succeeds, the program as a whole covers the union of the regimes.
\end{remark}
